\documentclass[10pt,a4paper]{amsart}
\usepackage[margin=19mm]{geometry}
\usepackage{amsmath,amssymb,mathtools}
\usepackage[scr=rsfs,cal=euler]{mathalfa}
\usepackage{xfrac}
\usepackage[unicode,hidelinks,bookmarks=false]{hyperref}
\newcommand{\ie}{i.e.\ }
\newcommand{\cf}{cf.\ }
\newcommand{\resp}{respectively\ }
\newcommand{\Subsection}{\subsection}
\providecommand{\nobreakdash}{\nobreak\hbox{-}}
\setlength{\emergencystretch}{2em}
\multlinegap0pt
\usepackage{amssymb}
\newtheorem{lemma}{Lemma}
\newtheorem{theorem}{Theorem}
\newtheorem{prop}{Proposition}
\newcommand{\SUM}[3]{\sum\limits_{{#1}={#2}}^{#3}}
\newcommand{\dif}{\mathrm{d}}
\title{Asymptotic stability of composite waves of two viscous shocks for relaxed compressible Navier-Stokes equations}
\author{Renyong Guan, Yuxi Hu}
\date{Source: arXiv:2408.17261v3 (CC BY 4.0)}
\begin{document}
\maketitle

\noindent\emph{Source and licence.} Asymptotic stability of composite waves of two viscous shocks for relaxed compressible Navier-Stokes equations. Version arXiv:2408.17261v3 (CC BY 4.0), distributed by arXiv under the Creative Commons Attribution 4.0 International licence, http://creativecommons.org/licenses/by/4.0/, as recorded in the arXiv licence metadata for this exact version and shown on the abstract page https://arxiv.org/abs/2408.17261v3. Full licence notice: \texttt{licenses/arxiv-2408.17261-CC-BY-4.0.txt}.\par\smallskip

\noindent\textit{Editorial scope.} Counted unit: the article's lemma Le4.7 (source0.tex lines 1541--1557). The article's complete original proof of it is delivered with it (source0.tex lines 1558--1842, one \texttt{proof} environment of 285 lines). No mathematics is written, repaired or re-derived: the statement and the proof are the article's own lines, in the article's own notation, and no retained line is substituted or re-flowed.  Retained uncounted as the source-local context the proof needs: lines 61--67; lines 75--77; lines 154--158; lines 272--289; lines 429--438; lines 479--506; lines 511--516; lines 541--579; lines 560--562; lines 772--780; lines 831--833; lines 843--845; lines 846--854.  Nothing was omitted from any retained span, so the package carries no omission record.  No retained line contains a citation, so the wrapper carries no bibliography and no bibliography entry.  No retained line contains a figure environment, an image-inclusion command or drawing code, so no image is delivered and the package declares no asset dependency.  Editorial formatting only, in addition: an AMS \texttt{amsart} preamble that restates the article's own declarations and macros used by the retained lines, namely the environments \texttt{lemma}, \texttt{theorem}, \texttt{prop} on separate counters, together with the standard packages those lines need.  The retained environments print in this document as Lemma 1, Theorem 1, Proposition 1 in order of appearance, whereas the article prints the same items with the numbering of its own full document.  The complete native source of the article as distributed in the arXiv e-print for this exact version is bundled unchanged as \texttt{sources/164511/source0.tex}.
\begin{equation}\label{1.3}
\begin{cases}
v_t-u_x=0,\\
u_t+p_x=\Pi_x,\\
\tau \Pi_t+v\Pi=\mu u_x,
\end{cases}
\end{equation}

\begin{align} \label{1.4}
(v, u, \Pi)(0,x)=(v_0, u_0, \Pi_0)(x)\rightarrow(v_{\pm},u_{\pm},0) \quad (x\rightarrow\pm\infty).
\end{align}

\begin{equation}\label{2.5}
\tau\leq \min\{\inf\limits_{z_1\in[v_m,v_-]}\frac{\mu}{2|\sigma_1^2+p^{\prime}(z_1)|},
\inf\limits_{z_2\in[v_m,v_+]}\frac{\mu}{2|\sigma_2^2+p^{\prime}(z_2)|},
1\}.
\end{equation}

\begin{lemma}\label{le2.1}
Let \eqref{2.5} hold.
 For any states $(v_-,u_-),(v_m,u_m),(v_+,u_+)\in \mathbb{R}_+\times\mathbb{R}$ and $\sigma_1<0,\sigma_2>0$ satisfying R-H condition and Lax condition, there exists a positive constant $C$ independent of $\tau$ such that the following is true: the traveling wave solutions $(\widetilde{u}_1,\widetilde{v}_1,\widetilde{\Pi}_1)(\xi_1)$ connecting $(v_-,u_-,0)$ and $(v_m,u_m,0)$ and $(\widetilde{u}_2,\widetilde{v}_2,\widetilde{\Pi}_2)(\xi_2)$ connecting $(v_m,u_m,0)$ and $(v_+,u_+,0)$ exist uniquely and satisfy
\begin{align*}
(\widetilde{v}_1)_{\xi_1}<0,\quad
(\widetilde{v}_2)_{\xi_2}>0,\quad
(\widetilde{v}_i)_{\xi_i}\sim(\widetilde{u}_i)_{\xi_i},
\end{align*}
and
\begin{align*}
&|\widetilde{v}_i(\xi_i)-v_m|\leq C\delta_i e^{-C\delta_i|\xi_i|},
|\widetilde{u}_i(\xi_i)-u_m|\leq C\delta_i e^{-C\delta_i|\xi_i|},\quad  \forall (-1)^i\xi_i<0,\\
&|\widetilde{\Pi}_i|\leq C\delta_i^2 e^{-C\delta_i|\xi_i|},\quad \quad
|\partial_{\xi_i}(\widetilde{v}_i,\widetilde{u}_i)|\leq C\delta_i^2e^{-C\delta_i|\xi_i|},  \,   |\partial_{\xi_i} \widetilde{\Pi}_i|\le C \delta_i |(\widetilde v_i)_{\xi_i}| \quad \forall\xi_i\in\mathbb{R},\\
&|\partial^k_{\xi_i}(\widetilde{v}_i,\widetilde{u}_i,\widetilde{\Pi}_i)|\leq C\delta_i|\partial_{\xi_i}\widetilde{v}_i|, \qquad \forall\xi_i\in\mathbb{R},
\end{align*}
for $i=1,2$ and $k=2,3,4$, where $\delta_i$ denote the strength of the shock as $\delta_1:=|p(v_-)-p(v_m)|\sim|v_--v_m|\sim|u_--u_m|$ and $\delta_2:=|p(v_m)-p(v_+)|\sim|v_m-v_+|\sim|u_m-u_+|$.
\end{lemma}

\begin{equation}\label{2.10}
\begin{cases}
\begin{aligned}
&\dot{X}_i(t)=-\frac{M}{\delta_i}
\left[\int_{\mathbb{R}}\frac{a}{\sigma_i}(\widetilde{u}_i^{X_i})_x(p(v)-p(\widetilde{v}))\dif x
-\int_{\mathbb{R}}a\left(p(\widetilde{v}_i^{X_i})\right)_x(v-\widetilde{v})\dif x\right],\\
&X_i(t)=0,
\end{aligned}
\end{cases}
\end{equation}

\begin{theorem}\label{th3.1}
Let $\underline{v}$ and $\underline{u}$ be smooth monotone functions such that
\[
\underline{v}(x)=v_{\pm} \quad \underline{u}(x)=u_{\pm} \quad for \quad \pm x\geq1.
\]
For any constants $M_0,M_1,\kappa_1,\kappa_2,\kappa_3,\kappa_4$ with $M_1>M_0>0$ and $\kappa_1>\kappa_2>\kappa_3>\kappa_4>0$, there exists a constant $T_0>0$ such that if
\[
\begin{aligned}
&\|v_0-\underline{v}\|_{H^2}+\|u_0-\underline{u}\|_{H^2}+\sqrt{\tau}\|\Pi_0\|_{H^2}\leq M_0,\\
&0<\kappa_3\leq v_0(x)\leq\kappa_2,\qquad \forall x\in\mathbb{R},
\end{aligned}
\]
then the system \eqref{1.3} with initial condition \eqref{1.4} has a unique solution $(v, u, \Pi)$ on $[0,T_0]$ such that
\[
\begin{aligned}
(v-\underline{v}, u-\underline{u}, \Pi)\in C([0,T_0];H^2),
\end{aligned}
\]
and
\[
\|v-\underline{v}\|_{L^{\infty}(0,T_0;H^2)}+\|u-\underline{u}\|_{L^{\infty}(0,T_0;H^2)}
+\sqrt{\tau}\|\Pi\|_{L^{\infty}(0,T_0;H^2)}\leq M_1.
\]
Moreover:
\[
\kappa_4\leq v(t,x)\leq\kappa_1,\qquad \forall(t,x)\in[0,T_0]\times\mathbb{R}.
\]
\end{theorem}

\begin{align}\label{3.1}
\left(\widetilde{v}, \widetilde{u}, \widetilde{\Pi}\right)(t, x):=
\left(\widetilde{v}_1^{X_1}+\widetilde{v}_2^{X_2}-v_m,
\widetilde{u}_1^{X_1}+\widetilde{u}_2^{X_2}-u_m,
\widetilde{\Pi}_1^{X_1}+\widetilde{\Pi}_2^{X_2}\right),
\end{align}

\begin{prop}\label{p1}
Let $(v, u, \Pi)$ be local solutions given by Theorem \ref{th3.1} on $[0,T]$ for some $T>0$ and $(\widetilde{v},\widetilde{u},\widetilde{\Pi})$ be defined in \eqref{3.1}.  Then there exist positive constants $\delta_0,\varepsilon_1$ such that if two independent shock waves strength satisfy $\delta_1,\delta_2<\delta_0$ and
\begin{equation}\label{3.5}
\sup_{0\le t\le T} \|(v-\widetilde{v}, u-\widetilde{u}, \sqrt{\tau}(\Pi-\widetilde{\Pi}))\|_{H^2}\leq\varepsilon_1,
\end{equation}
then the following estimates hold
\begin{equation}\label{3.6}
\begin{aligned}
&\sup\limits_{t\in[0,T]}\left(\|v-\widetilde{v}\|^2_{H^2}
+\|u-\widetilde{u}\|^2_{H^2}+\tau\|\Pi-\widetilde{\Pi}\|^2_{H^2}\right)
+\int_0^T\sum\limits_{i=1}\limits^2\delta_i|\dot{X}_i|^2\dif t\\
&\qquad\qquad+\int_0^T\left(G^s(U)
+\|\left((v-\widetilde{v})_{x}, (u-\widetilde{u})_{x}\right)\|^2_{H^1}+\|\Pi-\widetilde{\Pi}\|^2_{H^2}\right)\dif t\\
%+\int_0^T\left(\|S-\widetilde{S}_{-X}\|^2_{H^2(\mathbb{R})}\right)dt\\
&\leq C_0\left(\|v_0-\widetilde{v}_0(\cdot)\|^2_{H^2}+\|u_0-\widetilde{u}_0(\cdot)\|^2_{H^2}
+\tau\|\Pi_0-\widetilde{\Pi}_0(\cdot)\|^2_{H^2}\right)+C_0\delta_0.
\end{aligned}
\end{equation}
Here $C_0$ independent of $T$ and $\tau$ and
\begin{equation}\label{3.7}
G^s(U):=\sum\limits_{i=1}\limits^2\int_\mathbb{R}|(\widetilde{v}_i)^{X_i}_{x}| |\phi_i(v -\widetilde{v})|^2\dif x
\end{equation}
where cutoff functions $\phi_1,\phi_2$ are defined as follows
\begin{equation}\label{ctf9.15}
\begin{aligned}
&\phi_1:=
\begin{cases}
1,\qquad\qquad\qquad \qquad \qquad \qquad \qquad    if \quad  x<\frac{X_1(t)+\sigma_1t}{2},  \\
0,\qquad\qquad\qquad \qquad \qquad\qquad \qquad     if \quad  x>\frac{X_1(t)+\sigma_1t}{2},  \\
\text{linearly decreasing 1 to 0} \qquad\qquad   if \quad  \frac{X_1(t)+\sigma_1t}{2}\ge x\le \frac{X_1(t)+\sigma_1t}{2},
\end{cases}\\
&\phi_2(x):=1-\phi_1(x).
\end{aligned}
\end{equation}
In addition, by \eqref{2.10},
\begin{equation}\label{3.8}
|\dot{X}_1(t)|+|\dot{X}_2(t)|\leq C_0\|(v-\widetilde{v})(t,\cdot)\|_{L^{\infty}},\qquad\forall t\leq T.
\end{equation}
\end{prop}

\begin{equation}
G^s(U):=\sum\limits_{i=1}\limits^2\int_\mathbb{R}|(\widetilde{v}_i)^{X_i}_{x}| |\phi_i(v -\widetilde{v})|^2\dif x
\end{equation}

\begin{equation}\label{4.1}
\begin{cases}
(v-\widetilde{v})_t-\SUM i 1 2 \dot{X}_i \cdot (\widetilde{v}_i)^{X_i}_x-(u-\widetilde{u})_x=0,\\
(u-\widetilde{u})_t-\SUM i 1 2 \dot{X}_i \cdot (\widetilde{u}_i)^{X_i}_x+\left(p(v)-p(\widetilde{v})\right)_x
=(\Pi-\widetilde{\Pi})_x-F_1,\\
\tau(\Pi-\widetilde{\Pi})_t-\tau\SUM i 1 2\dot{X}_i \cdot (\widetilde{\Pi}_i)^{X_i}_x+v\Pi-\widetilde{v}\widetilde{\Pi}
=\mu(u-\widetilde{u})_x-F_2,
\end{cases}
\end{equation}

\begin{align}\label{26.1b}
 G(U):=\int_\mathbb{R}\frac{v}{\mu}|\Pi-\widetilde{\Pi}|^2\dif x,
\end{align}

\begin{equation}\label{3.10-1}
X_1(t)+\sigma_1t\le \frac{\sigma_1}{2}t,\qquad X_2(t)+\sigma_2t\ge \frac{\sigma_2}{2}t, \qquad t >0.
\end{equation}

\begin{lemma}\label{le4.3}
Assume \eqref{3.10-1}. Given $v_+>0$, there exist positive constants $\delta_0,C$ such that for any $\delta_1,\delta_2\in(0,\delta_0)$, the following estimates hold. For each $i=1,2$,
\begin{align*}
&|(\widetilde{v}_i)_x^{X_i}||\widetilde{v}-\widetilde{v}_i^{X_i}|\le C\delta_i\delta_1\delta_2e^{-C\min\{\delta_1,\delta_2\}t},\quad t>0, \quad x\in \mathbb{R},\\
&\int_\mathbb{R}|(\widetilde{v}_i)_x^{X_i}||\widetilde{v}-\widetilde{v}_i^{X_i}|\dif x\le C\delta_1\delta_2e^{-C\min\{\delta_1,\delta_2\}t},\quad t>0,\\
&\int_\mathbb{R}|(\widetilde{v}_1)_x^{X_1}||(\widetilde{v}_2)_x^{X_2}|\dif x\le
C\delta_1\delta_2e^{-C\min\{\delta_1,\delta_2\}t},\quad t>0.
\end{align*}
\end{lemma}

\begin{lemma}\label{Le4.7}
Under the hypotheses of Proposition \ref{p1}, there exists $C>0$ (independent of $\tau$ and $T$) such that for all $t\in(0,T]$, we have
\begin{equation}\label{26.2}
\begin{aligned}
&\|\left(\partial_{x}\Phi, \partial_{x}\Psi\right)\|_{H^1}^2
+\tau\|\partial_{x}Q\|_{H^1}^2
+\int_0^t\|\partial_{x}Q\|_{H^1}^2\dif t\leq
C(\|\left(\partial_{x}\Phi_0, \partial_{x}\Psi_0\right)\|_{H^1}^2
+\tau\|\partial_{x}Q_0\|_{H^1}^2)\\
 &+C\delta_0\int_0^tG^s(U)\dif t
+C(\delta_0+\varepsilon_1)\int_0^t\left(\|\left(\partial_{x}\Phi, \partial_{x}\Psi\right)\|_{H^1}^2+G(U)\right)\dif t
+C\int_0^t\sum\limits_{i=1}^2\delta_i^2|\dot{X}_i|^2\dif t
+C\delta_0,
\end{aligned}
\end{equation}
where $\Phi_0=v_0-\widetilde{v}_0(\xi), \Psi_0=u_0-\widetilde{u}_0(\xi), Q_0=\Pi_0-\widetilde{\Pi}_0(\xi)$, and $G^s(U), G(U)$ are defined in \eqref{3.7} and \eqref{26.1b}, respectively.
\end{lemma}

\begin{proof}
Applying $\partial^k_x(k=1,2)$ to the system \eqref{4.1}, we derive that
\begin{equation}\label{4.10}
\begin{cases}
\partial_{t}\partial_{x}^k\Phi
-\sum\limits_{i=1}\limits^2\dot{X}_i\partial_{x}^{k+1}\widetilde{v}_i^{X_i}-\partial_{x}^{k+1}\Psi=0,\\
\partial_{t}\partial_{x}^k\Psi
-\sum\limits_{i=1}\limits^2\dot{X}_i\partial_{x}^{k+1}\widetilde{u}_i^{X_i}
          +p^{\prime}(v)\partial_{x}^{k+1}\Phi=\partial_{x}^{k+1}Q+F_3^k,\\
\tau\partial_{t}\partial_{x}^kQ
-\tau\sum\limits_{i=1}\limits^2\dot{X}_i\partial_{x}^{k+1}\widetilde{\Pi}_i^{X_i}
             +v\partial_{x}^{k}Q=\mu\partial_{x}^{k+1}\Psi+F_4^k,
\end{cases}
\end{equation}
where $$F_3^k=p^{\prime}(v)\partial^{k+1}_{x}(v-\widetilde{v})-\partial^{k+1}_{x}(p(v)-p(\widetilde{v}))
-\partial^k_{x}F_1$$
and
$$F_4^k=v\partial^k_{x}(\Pi-\widetilde{\Pi})-\partial^k_{x}(v\Pi-\widetilde{v}\widetilde{\Pi})
-\partial^k_{x}F_2.$$

Multiplying the above equations by $-p^{\prime}(v)\partial_{x}^k\Phi,\partial_{x}^k\Psi,\frac{1}{\mu}\partial_{x}^kQ$, respectively, and integrating over $\mathbb{R}$, we get
\begin{align}\label{4.11}
  \frac{\dif}{\dif t}\int_\mathbb{R}\left(-\frac{p^{\prime}(v)}{2}(\partial_{x}^k\Phi)^2
  +\frac{1}{2}(\partial_{x}^k\Psi)^2
  +\frac{\tau}{2\mu}(\partial_{x}^kQ)^2\right)\dif x
   +\int_\mathbb{R}\frac{v}{\mu}(\partial_x^kQ)^2\dif x
    =:\sum\limits_{i=1}\limits^{8}R^k_i,
\end{align}
where
\[
R^k_1=-\int_\mathbb{R}\frac{p^{\prime\prime}(v)}{2}v_t(\partial_{x}^k\Phi)^2\dif x, \quad
R^k_2=\int_\mathbb{R}p^{\prime\prime}(v)v_{x}\partial_{x}^k\Phi\partial_{x}^k\Psi \dif x,
\]
\[
R^k_3=-\sum\limits_{i=1}\limits^2
\int_\mathbb{R}\dot{X}_i\partial_{x}^{k+1}\widetilde{v}^{X_i}_ip^{\prime}(v)\partial_{x}^k\Phi \dif x,
\quad
R^k_4=\sum\limits_{i=1}\limits^2
\int_\mathbb{R}\dot{X}_i\partial_{x}^{k+1}\widetilde{u}^{X_i}_i\partial_{x}^k\Psi \dif x,
\]
\[
R^k_5=\frac{\tau}{\mu}\sum\limits_{i=1}\limits^2
\int_\mathbb{R}\dot{X}_i\partial_{x}^{k+1}\widetilde{\Pi}^{X_i}_iQ\dif x,\quad
R^k_6=\int_\mathbb{R}F_3^k\partial_{x}^k\Psi \dif x,\quad
R^k_7=\frac{1}{\mu}\int_\mathbb{R}F_4^k\partial_{x}^kQ\dif x.
\]
Firstly, using Lemmas \ref{le2.1} and Sobolev's imbedding theorem, we have
\begin{align*}
R^k_1=-\int_\mathbb{R}\frac{p^{\prime\prime}(v)}{2}(u-\widetilde{u})_x(\partial_{x}^k\Phi)^2\dif x
-\int_\mathbb{R}\frac{p^{\prime\prime}(v)}{2}\widetilde{u}_x(\partial_{x}^k\Phi)^2\dif x
\le C(\varepsilon_1+\delta_1^2+\delta_2^2)\int_\mathbb{R}(\partial_{x}^k\Phi)^2\dif x.
\end{align*}
Similarly, for $R^k_2$, using Lemma \ref{le2.1}, Sobolev's imbedding theorem and Young's inequality, we have
\begin{align*}
R^k_2\le C(\varepsilon_1+\delta_1^2+\delta_2^2)\left(\int_\mathbb{R}(\partial_{x}^k\Phi)^2\dif x
+\int_\mathbb{R}(\partial_{x}^k\Psi)^2\dif x\right).
\end{align*}
For $R^k_3$, using Lemma \ref{le2.1} and Young's inequality, we have
\begin{align*}
\int_\mathbb{R}\dot{X}_i\partial_{x}^{k+1}\widetilde{v}^{X_i}_xp^{\prime}(v)\partial_{x}^k\Phi \dif x
&\le C\delta_i|\dot{X}_i|^2\int_\mathbb{R}|\partial_{x}^{k+1}\widetilde{v}^{X_i}_i|\dif x
+\frac{C}{\delta_i}\int_\mathbb{R}|\partial_{x}^{k+1}\widetilde{v}^{X_i}_i|(\partial_{x}^k\Phi)^2 \dif x\\
&\le C\delta_i^2|\dot{X}_i|^2
+C\delta_i\int_\mathbb{R}(\partial_{x}^k\Phi)^2 \dif x.
\end{align*}
So, we get
\[
R^k_3\le C\sum\limits_{i=1}\limits^2\left(\delta_i^2|\dot{X}_i|^2
+\delta_i\int_\mathbb{R}(\partial_{x}^k\Phi)^2 \dif x\right).
\]
Similarly, for $R^k_4,R^k_5$, we have
\[
R^k_4\le C\sum\limits_{i=1}\limits^2\left(\delta_i^2|\dot{X}_i|^2
+\delta_i\int_\mathbb{R}(\partial_{x}^k\Psi)^2 \dif x\right),\quad
R^k_5\le C\tau\sum\limits_{i=1}\limits^2\left(\delta_i^2|\dot{X}_i|^2
+\delta_i\int_\mathbb{R}(\partial_{x}^kQ)^2 \dif x\right).
\]
Next, we estimate $R_6^k$. For $k=1$, since
\begin{align*}
  &p^{\prime}(v)(v-\widetilde{v})_{xx}-(p(v)-p(\widetilde{v}))_{xx}\\
   &=-p^{\prime\prime}(v)v^2_{x}-p^{\prime}(v)\widetilde{v}_{xx}
                    +p^{\prime\prime}(\widetilde{v})\widetilde{v}^2_{x}
                    +p^{\prime}(\widetilde{v})\widetilde{v}_{xx}\\
   &=-p^{\prime\prime}(v)(v_{x}-\widetilde{v}_{x})^2
   -2p^{\prime\prime}(v)(v_{x}-\widetilde{v}_{x})\widetilde{v}_{x}
     -(p^{\prime\prime}(v)-p^{\prime\prime}(\widetilde{v}))\widetilde{v}^2_{x}
     -(p^{\prime}(v)-p^{\prime}(\widetilde{v}))\widetilde{v}_{xx},
\end{align*}
using Lemma \ref{le2.1}, we have
\begin{align*}
 &\int_\mathbb{R}(u-\widetilde{u})_{x}(p^{\prime}(v)(v-\widetilde{v})_{xx}-(p(v)-p(\widetilde{v}))_{xx})\dif x\\
      &\leq C\int_\mathbb{R}|(u-\widetilde{u})_{x}||(v-\widetilde{v})_{x}|^2\dif x
      +C\int_\mathbb{R}|(u-\widetilde{u})_{x}||(v-\widetilde{v})_{x}||\widetilde{v}_{x}|\dif x
      +C\int_\mathbb{R}|(u-\widetilde{u})_{x}||v-\widetilde{v}||\widetilde{v}_{x}|\dif x.
\end{align*}
We estimate each term of the right hand side of the above inequality. Using Young's inequality and Sobolev's imbedding theorem, we have
\[
\int_\mathbb{R}|(u-\widetilde{u})_{x}||(v-\widetilde{v})_{x}|^2\dif x\leq
 C\varepsilon_1\int_\mathbb{R}|\partial_{x}\Phi|^2\dif x
\]
and
\begin{align*}
\int_\mathbb{R}|(u-\widetilde{u})_{x}||(v-\widetilde{v})_{x}||\widetilde{v}_{x}|\dif x\leq C(\delta_1^2+\delta_2^2)
(\int_\mathbb{R}|\partial_{x}\Phi|^2\dif x+\int_\mathbb{R}|\partial_{x}\Psi|^2\dif x).
\end{align*}
In a similar way, using Lemma \ref{le2.1} and Young's inequality, we get
\begin{align*}
&\int_\mathbb{R}|(u-\widetilde{u})_{x}||v-\widetilde{v}||(\widetilde{v}_i)^{X_i}_{x}|\dif x\\
&\le C\delta_i\left(\int_\mathbb{R}|\partial_{x}\Psi|^2\dif x
+\int_\mathbb{R}|(\widetilde{v}_i)^{X_i}_{x}||v-\widetilde{v}|^2\dif x\right)\\
&\le  C\delta_i\left(\int_\mathbb{R}|\partial_{x}\Psi|^2\dif x
+\int_\mathbb{R}|(\widetilde{v}_i)^{X_i}_{x}||\phi_i(v-\widetilde{v})|^2\dif x
+\int_\mathbb{R}|(\widetilde{v}_i)^{X_i}_{x}|(1-\phi_i^2)|v-\widetilde{v}|^2\dif x\right)\\
&\le C\delta_i\left(\int_\mathbb{R}|\partial_{x}\Psi|^2\dif x
+G^s(U)
+\varepsilon_1 \delta_i^2e^{-C\delta_i t}\right).
\end{align*}
Hence, we derive that
\begin{align*}
 &\int_\mathbb{R}(u-\widetilde{u})_{x}(p^{\prime}(v)(v-\widetilde{v})_{xx}-(p(v)-p(\widetilde{v}))_{xx})\dif x\\
  &\le C(\varepsilon_1+\delta_1+\delta_2)\int_\mathbb{R}|\partial_{x}\Phi|^2\dif x
  +C(\delta_1+\delta_2)\int_\mathbb{R}|\partial_{x}\Psi|^2\dif x
  +C(\delta_1+\delta_2)G^s+C\sum\limits_{i=1}\limits^2\varepsilon_1 \delta_i^3e^{-C\delta_i t}.
\end{align*}
On the other hand, using Lemma \ref{le2.1}, we derive that for $k=1,2$,
\begin{align*}
\partial_x^{k}F_1&=
\partial_x^{k+1}\left(p(\widetilde{v})-p(\widetilde{v}_1^{X_1})-p(\widetilde{v}_2^{X_2})\right)\\
&\le
C\left(|(\widetilde{v}_1)^{X_1}_x||\widetilde{v}-(\widetilde{v}_1)^{X_1}|
+|(\widetilde{v}_2)^{X_2}_x||\widetilde{v}-(\widetilde{v}_2)^{X_2}|
+|(\widetilde{v}_1)^{X_1}_x||(\widetilde{v}_2)^{X_2}_x|\right).
\end{align*}
So, using Lemmas \ref{le2.1}, \ref{le4.3}, H\"{o}lder inequality and Sobolev's imbedding theorem, we have
\begin{align*}
-&\int_{\mathbb{R}}(u-\widetilde{u})_{x}\partial_{x}F_1\dif x\\
&\le C\varepsilon_1 \left(\left\||(\widetilde{v}_1)^{X_1}_x||\widetilde{v}-(\widetilde{v}_1)^{X_1}|\right\|_{L^2}
+\left\||(\widetilde{v}_2)^{X_2}_x||\widetilde{v}-(\widetilde{v}_2)^{X_2}|\right\|_{L^2}
+\left\||(\widetilde{v}_1)^{X_1}_x||(\widetilde{v}_2)^{X_2}_x|\right\|_{L^2}\right)\\
&\le C\varepsilon_1\delta_1\delta_2(\delta_1^{\frac{1}{2}}+\delta_2^{\frac{1}{2}})
e^{-C\min\{\delta_1,\delta_2\}t}.
\end{align*}
Combining the above estimates, we derive that
\begin{align*}
R_6^1\le &C(\varepsilon_1+\delta_1+\delta_2)\int_\mathbb{R}|\partial_{x}\Phi|^2\dif x
  +C(\delta_1+\delta_2)\int_\mathbb{R}|\partial_{x}\Psi|^2\dif x
  +C(\delta_1+\delta_2)G^s\\
  &+C\varepsilon_1\delta_1\delta_2e^{-C\min\{\delta_1,\delta_2\}t}+C\sum\limits_{i=1}\limits^2\varepsilon_1 \delta_i^3e^{-C\delta_i t}.
\end{align*}
Similarly, for $k=2$, since
\begin{align*}
  &p^{\prime}(v)(v-\widetilde{v})_{xxx}-(p(v)-p(\widetilde{v}))_{xxx}\\
   &=-p^{\prime\prime\prime}(v)(v_{x}-\widetilde{v}_{x})^3
   -(p^{\prime\prime\prime}(v)-p^{\prime\prime\prime}(\widetilde{v}))\widetilde{v}_{x}^3
   -3p^{\prime\prime\prime}(v)\widetilde{v}_{x}^2(v_{\xi}-\widetilde{v}_{x})
   -3p^{\prime\prime\prime}(v)\widetilde{v}_{x}(v_{x}-\widetilde{v}_{x})^2\\
   &\quad-3p^{\prime\prime}(v)(v_{x}-\widetilde{v}_{x})(v_{xx}
   -\widetilde{v}_{xx})-3p^{\prime\prime}(v)\widetilde{v}_{x}(v_{xx}-\widetilde{v}_{xx})
  -3p^{\prime\prime}(v)\widetilde{v}_{xx}(v_{x}-\widetilde{v}_{x})\\
   &\quad-3(p^{\prime\prime}(v)-p^{\prime\prime}(\widetilde{v}))\widetilde{v}_{x}\widetilde{v}_{xx}
   -(p^{\prime}(v)-p^{\prime}(\widetilde{v}))\widetilde{v}_{xxx}.
\end{align*}
Using Lemma \ref{le2.1}, Young's inequality and Sobolev's imbedding theorem, we have
\begin{align*}
 &\int_\mathbb{R}(u-\widetilde{u})_{xx}
     (p^{\prime}(v)(v-\widetilde{v})_{xxx}-(p(v)-p(\widetilde{v}))_{xxx})\dif x\\
      &\leq C\int_\mathbb{R}|(u-\widetilde{u})_{xx}||v-\widetilde{v}||\widetilde{v}_x|\dif x
      +C\int_\mathbb{R}|(u-\widetilde{u})_{xx}||(v-\widetilde{v})_{x}|^3\dif x\\
     &\quad +C\int_\mathbb{R}|(u-\widetilde{u})_{xx}||(v-\widetilde{v})_x||\widetilde{v}_x|\dif x
      +C\int_\mathbb{R}|(u-\widetilde{u})_{xx}||(v-\widetilde{v})_{xx}||\widetilde{v}_x|\dif x\\
& \le C(\varepsilon_1^2+\delta_1+\delta_2)\left(\int_\mathbb{R}|\partial_{x}\Phi|^2\dif x
+\int_\mathbb{R}|\partial_{xx}\Psi|^2\dif x\right)
  +C(\delta_1+\delta_2)\int_\mathbb{R}|\partial_{xx}\Phi|^2\dif x\\
 &\quad +C(\delta_1+\delta_2)G^s+C\sum\limits_{i=1}\limits^2\varepsilon_1 \delta_i^3e^{-C\delta_i t}
\end{align*}
and
\begin{align*}
-&\int_{\mathbb{R}}(u-\widetilde{u})_{xx}\partial_{xx}F_1\dif x\\
&\le C\varepsilon_1 \left(\left\||(\widetilde{v}_1)^{X_1}_x||\widetilde{v}-(\widetilde{v}_1)^{X_1}_x|\right\|_{L^2}
+\left\||(\widetilde{v}_2)^{X_2}_x||\widetilde{v}-(\widetilde{v}_2)^{X_2}_x|\right\|_{L^2}
+\left\||(\widetilde{v}_1)^{X_1}_x||(\widetilde{v}_2)^{X_2}_x|\right\|_{L^2}\right)\\
&\le C\varepsilon_1\delta_1\delta_2(\delta_1^{\frac{1}{2}}+\delta_2^{\frac{1}{2}})
e^{-C\min\{\delta_1,\delta_2\}t}.
\end{align*}
So, we have
\begin{align*}
&R_6^2\le C(\varepsilon_1^2+\delta_1+\delta_2)\left(\int_\mathbb{R}|\partial_{x}\Phi|^2\dif x
+\int_\mathbb{R}|\partial_{xx}\Psi|^2\dif x\right)
  +C(\delta_1+\delta_2)\int_\mathbb{R}|\partial_{xx}\Phi|^2\dif x\\
   &\qquad+C(\delta_1+\delta_2)G^s
  +C\varepsilon_1\delta_1\delta_2e^{-C\min\{\delta_1,\delta_2\}t}+C\sum\limits_{i=1}\limits^2\varepsilon_1 \delta_i^3e^{-C\delta_i t}.
\end{align*}
Similarly, for $R_7^1$, since
\[
v(\Pi-\widetilde{\Pi})_{x}-(v\Pi-\widetilde{v}\widetilde{\Pi})_{x}=-(v-\widetilde{v})\widetilde{\Pi}_{x}
-(v-\widetilde{v})_{x}(\Pi-\widetilde{\Pi})
-(v-\widetilde{v})_{x}\widetilde{\Pi}-\widetilde{v}_{x}(\Pi-\widetilde{\Pi}),
\]
using Lemma \ref{le2.1} and $|\widetilde{\Pi}_i^{X_i}|\le C|(\widetilde{v}_i)^{X_i}_x|$, we have
\begin{align*}
 &\int_\mathbb{R}\frac{(\Pi-\widetilde{\Pi})_{x}}{\mu}
 \left(v(\Pi-\widetilde{\Pi})_{x}-(v\Pi-\widetilde{v}\widetilde{\Pi})_{x}\right)\dif x\\
 &\le C\int_\mathbb{R}\frac{|(\Pi-\widetilde{\Pi})_{x}|}{\mu}|v-\widetilde{v}||\widetilde{v}_{x}|\dif x
 + C\int_\mathbb{R}\frac{|(\Pi-\widetilde{\Pi})_{x}|}{\mu}|(v-\widetilde{v})_{x}||\Pi-\widetilde{\Pi}|\dif x\\
 &\quad+ C\int_\mathbb{R}\frac{|(\Pi-\widetilde{\Pi})_{x}|}{\mu}|(v-\widetilde{v})_{x}||\widetilde{\Pi}|\dif x
 + C\int_\mathbb{R}\frac{|(\Pi-\widetilde{\Pi})_{x}|}{\mu}|\Pi-\widetilde{\Pi}||\widetilde{v}_x|\dif x\\
 &\le C(\varepsilon_1+\delta_1+\delta_2)\left(\int_\mathbb{R}\frac{v}{\mu}(\partial_xQ)^2\dif x+G\right)
 +C(\delta_1+\delta_2)\int_\mathbb{R}(\partial_x\Phi)^2\dif x
 +C(\delta_1+\delta_2)G^s
 +C\sum\limits_{i=1}\limits^2\varepsilon_1 \delta_i^3e^{-C\delta_i t}.
\end{align*}
On the other hand, for $k=1,2$, using Lemma \ref{le2.1}, we get
\begin{align*}
\partial_x^{k}F_2&=
\partial_x^{k}\left(\left(\widetilde{v}^{X_2}_2-v_m\right)\widetilde{\Pi}_1^{X_1}
+\left(\widetilde{v}^{X_1}_1-v_m\right)\widetilde{\Pi}_2^{X_2}\right)\\
&\le
C\left(|(\widetilde{v}_1)^{X_1}_x||\widetilde{v}-(\widetilde{v}_1)^{X_1}_x|
+|(\widetilde{v}_2)^{X_2}_x||\widetilde{v}-(\widetilde{v}_2)^{X_2}_x|
+|(\widetilde{v}_1)^{X_1}_x||(\widetilde{v}_2)^{X_2}_x|\right).
\end{align*}
Then, using Lemma \ref{le4.3} and Young's inequality, we have
\begin{align*}
 &\int_\mathbb{R}\frac{(\Pi-\widetilde{\Pi})_{x}}{\mu}
 \left(\left(\widetilde{v}^{X_2}_2-v_m\right)\widetilde{\Pi}_1^{X_1}
+\left(\widetilde{v}^{X_1}_1-v_m\right)\widetilde{\Pi}_2^{X_2}\right)_x\dif x\\
&\le C\|\partial_x F_2\|_{L^2} +\frac{1}{16}\int_\mathbb{R}\frac{v}{\mu}(\partial_xQ)^2\dif x\\
&\le C\delta_1^2\delta_2^2e^{-C\min\{\delta_1,\delta_2\}t}
+\frac{1}{16}\int_\mathbb{R}\frac{v}{\mu}(\partial_xQ)^2\dif x.
\end{align*}
Combining the above estimates, we derive that
\begin{align*}
 &R_7^1 \le \frac{1}{8}\int_\mathbb{R}\frac{v}{\mu}(\partial_xQ)^2\dif x
 +C(\varepsilon_1+\delta_1+\delta_2)G
 +C(\delta_1+\delta_2)\left(\int_\mathbb{R}(\partial_x\Phi)^2\dif x+G^s\right)\\
 &+C\delta_1^2\delta_2^2e^{-C\min\{\delta_1,\delta_2\}t}
 +C\sum\limits_{i=1}\limits^2\varepsilon_1 \delta_i^3e^{-C\delta_i t}.
\end{align*}
Next, for $R_7^2$, first we have
\begin{align*}
v(\Pi-\widetilde{\Pi})_{xx}-(v\Pi-\widetilde{v}\widetilde{\Pi})_{xx}=
&-(v-\widetilde{v})_{xx}(\Pi-\widetilde{\Pi})
-(v-\widetilde{v})_{xx}\widetilde{\Pi}-\widetilde{v}_{xx}(\Pi-\widetilde{\Pi})
-(v-\widetilde{v})\widetilde{\Pi}_{xx}\\
&-2(v-\widetilde{v})_{x}(\Pi-\widetilde{\Pi})_{x}-2(v-\widetilde{v})_{x}\widetilde{\Pi}_{x}
-2\widetilde{v}_{x}(\Pi-\widetilde{\Pi})_{x}.
\end{align*}
Specially, using H\"{o}lder inequality, Sobolev's imbedding theorem and Young's inequality, we have
\begin{align*}
&\int_\mathbb{R}(v-\widetilde{v})_{xx}(\Pi-\widetilde{\Pi})(\Pi-\widetilde{\Pi})_{xx}\dif x\\
&\leq C
\|(v-\widetilde{v})_{\xi\xi}\|_{L^2}
\left(\int_\mathbb{R}\left(\frac{v}{\mu}\right)^2(\Pi-\widetilde{\Pi})^2(\Pi-\widetilde{\Pi})_{xx}^2\dif x\right)^{\frac{1}{2}}\\
&\leq C\varepsilon_1|\sqrt{\frac{v}{\mu}}(\Pi-\widetilde{\Pi})|_{L^{\infty}}
\|\sqrt{\frac{v}{\mu}}(\Pi-\widetilde{\Pi})_{xx}\|_{L^2}\\
&\leq C\varepsilon_1\left(\|\sqrt{\frac{v}{\mu}}(\Pi-\widetilde{\Pi})\|_{H^1}^2+
\|\sqrt{\frac{v}{\mu}}(\Pi-\widetilde{\Pi})_{xx}\|_{L^2}^2\right).
\end{align*}
So, we have
\begin{align*}
 &\int_\mathbb{R}\frac{(\Pi-\widetilde{\Pi})_{xx}}{\mu}
 \left(v(\Pi-\widetilde{\Pi})_{xx}-(v\Pi-\widetilde{v}\widetilde{\Pi})_{xx}\right)\dif x\\
 &\le C(\varepsilon_1+\delta_1+\delta_2)\left\|\sqrt{\frac{v}{\mu}}Q\right\|_{H^2}
 +(\delta_1+\delta_2)\left\|\partial_x\Phi\right\|_{H^1}+(\delta_1+\delta_2)G^s+C\sum\limits_{i=1}\limits^2\varepsilon_1 \delta_i^3e^{-C\delta_i t}.
\end{align*}
On the other hand, using the definition of $F_2$ and Young's inequality, we have
\begin{align*}
 &\int_\mathbb{R}\frac{(\Pi-\widetilde{\Pi})_{xx}}{\mu}
 \left(\left(\widetilde{v}^{X_2}_2-v_m\right)\widetilde{S}_1^{X_1}
+\left(\widetilde{v}^{X_1}_1-v_m\right)\widetilde{\Pi}_2^{X_2}\right)_{xx}\dif x\\
&\le C\|\partial_{xx}F_2\|_{L^2}
+\frac{1}{16}\int_\mathbb{R}\frac{v}{\mu}(\partial_{xx}Q)^2\dif x\\
&\le C\delta_1^2\delta_2^2e^{-C\min\{\delta_1,\delta_2\}t}
+\frac{1}{16}\int_\mathbb{R}\frac{v}{\mu}(\partial_{xx}Q)^2\dif x.
\end{align*}
Therefore, we get
\begin{align*}
R_7^2&\le \frac{1}{8}\int_\mathbb{R}\frac{v}{\mu}(\partial_{xx}Q)^2\dif x
 +C(\varepsilon_1+\delta_1+\delta_2)\left\|\sqrt{\frac{v}{\mu}}Q\right\|_{H^1}
  +C(\delta_1+\delta_2)\left(G^s+\left\|\partial_x\Phi\right\|_{H^1}\right)\\
  &+C\sum\limits_{i=1}\limits^2\varepsilon_1 \delta_i^3e^{-C\delta_i t}
 +C\delta_1^2\delta_2^2e^{-C\min\{\delta_1,\delta_2\}t}.
\end{align*}
Finally, integrating the equality \eqref{4.11} over $[0,t]$, using the above estimates, we complete the proof of this lemma.
\end{proof}

\end{document}
